
\textbf{Proposition 3.4.} \uline{Avec les hypothèses et les notations de} 3.1 \uline{et} 3.3, \uline{on a pour tout} @@M0@@ \uline{la relation}:
@@M1@@

Comme les deux membres de la relation à démontrer sont additifs, on peut supposer que @@M2@@ est la classe d'un @@M3@@-Module localement libre de type fini @@M4@@. L'élément @@M5@@ est alors la classe du complexe @@M6@@, obtenu en prenant une résolution projective de @@M7@@ et en lui appliquant le foncteur @@M8@@. Les faisceaux d'homologie du complexe ainsi obtenu sont les @@M9@@. Or d'après 2.4 et 3.2, on a:
@@M10@@

Il en résulte que ces faisceaux d'homologie sont localement libres de type fini sur @@M11@@, donc parfaits sur @@M12@@; par suite, la classe dans @@M13@@ du complexe @@M14@@ est la somme alternée des classes de ses faisceaux d'homologie, soit:
@@M15@@

\textbf{Lemme 3.5.} \uline{Avec les hypothèses et les notations de} 3.1, \uline{on a les isomorphismes suivants} :
\begin{enumerate}[label=\roman*)]
\item \uline{pour tout} @@M16@@, \uline{et pour tout} @@M17@@, @@M18@@;
\item \uline{pour tout} @@M19@@, @@M20@@.
\end{enumerate}

On peut évidemment supposer @@M21@@ affine, donc quasi-compact.

i) Rappelons d'abord le lemme suivant, qui a été démontré dans exp. III.

